\documentclass{article}
\usepackage[left=2cm, right=2cm, top=2cm, bottom=2cm]{geometry}
\usepackage{graphicx}
\usepackage{amsmath}
\usepackage{amssymb}
\usepackage{amsthm}
\usepackage{fancyhdr}
\usepackage{verbatim}
\usepackage{listings}
\usepackage{xcolor}
\usepackage{pdfpages}
\usepackage{cancel}
\usepackage{physics}

\lstset{
    language=C++,
    basicstyle=\ttfamily\footnotesize,
    keywordstyle=\color{blue},
    commentstyle=\color{green},
    stringstyle=\color{red},
    numbers=left,
    numberstyle=\tiny\color{gray},
    frame=single,
    breaklines=true,
    captionpos=b,
}


\title{MTH 553: Lecture \# 20}
\author{Cliff Sun}

\newtheorem{theorem}{Theorem}[section]
\newtheorem{lemma}[theorem]{Lemma}
\newtheorem{definition}[theorem]{Definition}
\newtheorem{conjecture}[theorem]{Conjecture}
\newtheorem{proposition}[theorem]{Proposition}
\newtheorem{corollary}[theorem]{Corollary}
\newtheorem{one minute paper}[theorem]{One Minute Paper}

\pagestyle{fancy}
\lhead{\textbf{\thepage}\ \ \nouppercase{\rightmark}}
\chead{MTH 553: Lecture \# 20}
\rhead{Cliff Sun}

\begin{document}

\maketitle

\section*{Lecture Span}
\begin{enumerate}
    \item Fourier inversion
    \item Solving heat equation with Fourier transform
\end{enumerate}

For $g \in S$ (remember that $S$ is the Schwartz function). Then $\check{\hat{g}} = g$. Remark, Fourier transform and inversion also hold for $g \in L^2(\mathbb{R}^n)$. The Fourier transform is an 
\underline{$L^2$-isometry} 

\begin{equation}
    ||g||_{L^2} = ||\hat{g}||_{L^2}
\end{equation}


\section*{Heat kernel on $\mathbb{R}^n$ by Fourier transform}

Formally, $u_t = \triangle u$ and therefore 

\begin{equation}
    \hat{(u_t)} = \hat{\triangle u}
\end{equation}

\begin{align*}
    \implies \partial_t \hat{u}(\xi, t) &= \sum_{j=1}^n \left( \hat{\partial^2_x u} \right) \\
    &= \sum (i\xi_j)^{2}\hat{u}(\xi,t) \\
    &= i|\xi|^2 \hat{u}
\end{align*}

This is simply an ODE with respect to the $t$ coordinate. The solution is then 

\begin{equation}
    \hat{u}(\xi,t) = C(\xi)e^{-|\xi|^2 t}
\end{equation}

We can also use the initial condition that $u(x,0) = g(x)$. Then when $t=0$, we obtain that 

\begin{equation}
    \hat{u}(\xi,0) = \hat{g}(\xi)
\end{equation}

Therefore, 

\begin{equation}
    \hat{u}(\xi,t) = \hat{g}(\xi)e^{-|\xi|^2 t}
\end{equation}

Now, 

\begin{align*}
    u(x,t) &= \left[\hat{g}(\xi)e^{-|\xi|^2 t}\right]\check{} \\
    &= \frac{1}{(2\pi)^{n/2}}\int_{\mathbb{R}^n}e^{i \xi \cdot x}\hat{g}(\xi)e^{-|\xi|^2 t}d\xi\\
    &= \int_{\mathbb{R}^n}\left[ \frac{1}{(2\pi)^n} \int_{\mathbb{R}^n}e^{i(x-y) \dot \xi - |\xi|^2 t}d\xi \right]g(y)dy
\end{align*}

This last line was determined by the definition of $\hat{g}(x)$ using a dummy variable $y$. So the goal now is to evaluate this integral kernel. 

\subsection*{n=1}

Then this integral is  

\begin{align*}
    &\frac{1}{2\pi}\int_{\mathbb{R}}e^{i(x-y)\xi - \xi^2 t}d\xi \\
    &= \frac{e^{-(x-y)^2/4t}}{2\pi}\int_{\mathbb{R}^n}e^{-\left[ \xi - i\frac{(x-y)}{2t} \right]^2t}d\xi\\
    &= \frac{e^{-\left( x-y \right)^2/4t}}{2\pi}\int_{\gamma} e^{-z^2 t}dz
\end{align*}

Where $\gamma$ is a contour on the complex plane. Then, from complex analysis, we use Cauchy's theorem, and then we obtain the following integral:

\begin{align*}
    \frac{e^{-\left( x-y \right)^2/4t}}{2\pi}\int_{\gamma} e^{-z^2 t}dz &= \frac{e^{-\left( x-y \right)^2/4t}}{2\pi}\int_{\mathbb{R}^n} e^{-z^2 t}dz \\
    &= \frac{e^{-\left( x-y \right)^2/4t}}{2\pi}\frac{\sqrt{\pi}}{\sqrt{t}}
\end{align*}

Therefore, the heat kernel in 1-d is 

\begin{equation}
    \frac{1}{\sqrt{4\pi t}}e^{-(x-y)^2/4t}
\end{equation}


For $n \geq 1$, we can split up the integrals and splitting the exponential, getting 

\begin{equation}
    \frac{1}{(4\pi t)^{n/2}}e^{-|x-y|^2/4t} = K(x,y,t)
\end{equation}

This is the heat kernel. Now, what does this suggest? This implies a solution formula for the initial value problem. If $g(x)$ is bounded and continuous, then:

\begin{align*}
    u(x,t) &= \int_{\mathbb{R}^n}K(x,y,t)g(y)dy \\
    &= \frac{1}{(4\pi t)^{n/2}}\int_{\mathbb{R}^n}e^{-|x-y|^2/4t} g(y)dy
\end{align*}

Where $K(x,y,t)$ is defined as above. In fact, this is a bounded and $C^{\infty}$ smooth solution of 

\begin{equation}
    u_t = \triangle u \quad x \in \mathbb{R}^n, t > 0
\end{equation}

A few remarks:

\begin{enumerate}
    \item There exists an analogous heat kernel formula for bounded domains in Section 5.1 using eigenfunctions. 
    \item If $g$ is bounded but is only known to be continuous at point $\tilde{x}$, then the solution $u$ is continuous at $\tilde{x}$ (not for all $x$)
    \item Infinite propagation speed for heat equation. Suppose $g \geq 0$ with compact support and $g> 0$ somewhere. Then the heat equation states that $u(x,t)> 0$ is positive for all $x$ and $t$, regardless of 
    how small $t$ may be. Clearly \underline{non-physical phenomenon}. 
    \item Instanenous smoothing, i.e. $u(x,t)$ is $C^{\infty}$-smooth $\forall t > 0$ even though $g(x)$ is bounded (doesn't have to be continuous!) 
\end{enumerate}

\end{document}